%% The introduction is unnumbered by FI convention, so the heading and
%% table of contents entry are added manually.
\chapter*{Introduction}
\markright{\textsc{Introduction}}
\addcontentsline{toc}{chapter}{Introduction}

Security certifications are how buyers of IT products, particularly
in the public sector, are given assurance that a product meets a defined
set of security requirements. The two schemes studied in this thesis are
Common Criteria (CC)~\cite{iso15408}, which evaluates IT products
against a stated set of security claims, and FIPS~140~\cite{fips1402,fips1403},
which validates cryptographic modules. Both produce public documents, and
both describe products that are rarely built from scratch: cryptography
in particular tends to come from a small number of shared libraries.

OpenSSL is one of the most prominent of these libraries. The effects of
its widespread reuse became visible in 2014, when the Heartbleed
vulnerability (CVE-2014-0160) prompted large-scale measurement of how
many services were exposed and how quickly they were
patched~\cite{durumeric2014heartbleed}. In the same year, LibreSSL was
forked from OpenSSL~\cite{libressl} and Google announced
BoringSSL~\cite{langley2014boringssl}. Today, the library that a product
calls ``OpenSSL'' may be one of several related code bases.

\section*{Problem statement}

Certification documents are not designed to answer supply-chain questions.
Which third-party cryptographic library a product uses, and in which
version, is usually disclosed, if at all, in the prose of a PDF
document. The sec-certs project~\cite{janovsky2024seccerts} has made
thousands of these documents accessible and machine-readable, but tracing
the lineage of one specific library across Security Targets, certification
reports, and maintenance reports still requires a targeted analytical
pipeline. As a result, it is difficult to establish how far OpenSSL and
its forks have spread through certified products, or how far certified
products trail the upstream releases they depend on.

\section*{Goals and research questions}

This thesis investigates the footprint, usage patterns, and lifecycle of
OpenSSL and its forks in the certified products indexed by sec-certs.org,
and implements an extensible, reproducible framework for analyzing and
visualizing these dependencies. It addresses three research questions:

\begin{description}
  \item[RQ1] How widely are OpenSSL and its clones and forks represented
    across Common Criteria and FIPS~140 certifications indexed by
    sec-certs.org?
  \item[RQ2] What is the gap between upstream OpenSSL releases and the
    validation or re-validation of dependent certified products?
  \item[RQ3] How can automated data mining, AI/NLP document parsing, and
    interactive visual tooling be integrated into sec-certs workflows to
    reliably track cryptographic dependencies and produce reproducible
    business intelligence?
\end{description}

The thesis is an empirical study of public certification data. Manual
reverse engineering, side-channel analysis, and binary-level vulnerability
research of certified products are out of scope, as are certification
frameworks other than CC and FIPS~140.

\section*{Contributions}

The thesis makes two contributions:

\begin{itemize}
  \item an empirical analysis of vendor adoption trends, re-certification
    patterns, and compliance lag in open-source cryptography; and
  \item an open toolkit of Jupyter notebooks and processing scripts that
    extends sec-certs and supports ongoing monitoring and reproducible
    analysis.
\end{itemize}

\section*{Thesis outline}

Chapter~\ref{ch:background} introduces the certification schemes, the
OpenSSL family of libraries, the sec-certs project, and related work.
Chapter~\ref{ch:method} defines the dataset, the detection heuristics, the
AI-assisted parsing, and the metrics. Chapter~\ref{ch:results} presents the
findings for each research question, and Chapter~\ref{ch:discussion}
interprets them and considers threats to validity. The Conclusion summarizes
the contributions and proposes future work.
